\section{A finite resolution that computes cohomology}
\label{sec:resolution}

The geometry above serves two purposes. It bounds the size of the
space, and it provides a cover with projective free sheaves on all its
intersections. We spell out the latter point so that the passage from
functions to ordinary sheaf cohomology is explicit.

\subsection{Free sheaves on opens}

For an open \(W\) in a space \(X\), let \(P(W)\) be the free
abelian sheaf represented by \(W\). Concretely, sheafify the free
abelian presheaf on the representable presheaf of \(W\). What we
use is its natural representing property:

\[
\Hom(P(W),\mathcal M)\cong\Gamma(W,\mathcal M).
\]

An inclusion \(V\subseteq W\) induces \(P(V)\to P(W)\),
and \(P(X)\cong\underline{\Z}_X\). For a commutative ring \(R\),
write \(P_R(W)\) for the corresponding free sheaf of \(R\)-modules.
It represents sections in that category, with
\(P_R(X)\cong\underline R_X\).

\begin{lemma}\label{lem:projective-opens}
Suppose \(X\) is Hausdorff and totally disconnected. If \(W\) is
a disjoint union of compact-open subsets of \(X\), then \(P(W)\)
and \(P_R(W)\) are projective.\leanmark{projectivity}
\end{lemma}

\begin{proof}
To prove projectivity, it suffices to lift a section over \(W\)
through an epimorphism of sheaves. Epimorphisms are locally surjective,
so local lifts exist. On a compact open \(C\), choose finitely many
clopen sets refining a local lifting cover. This is possible because
\(C\) is compact Hausdorff and totally disconnected, hence
zero-dimensional. Subtract earlier clopen sets from later ones to
make the finite cover disjoint. The local lifts then glue, with no
compatibility conditions on overlaps, to a section on \(C\).
Do this on each compact-open component of \(W\) and glue again.
The representing property gives projectivity in both coefficient
categories.
\end{proof}

\subsection{The ordered complex and its exactness}

For the cover \(W_0,\ldots,W_r\), set

\[
P_k=\bigoplus_{0\leq i_0<\cdots<i_k\leq r}
P(W_{i_0}\cap\cdots\cap W_{i_k}).
\]

Deleting one index enlarges the intersection and induces a map of free
sheaves. Let the differential be the alternating sum of these maps.
The augmentation \(P_0\to\underline\Z_{U_r}\) is induced by the
inclusions of the charts. The usual face-deletion identities imply
that the differential squares to zero.

\begin{proposition}\label{prop:resolution}
The augmented complex

\[
0\longrightarrow P_r\longrightarrow\cdots\longrightarrow P_1
\longrightarrow P_0\longrightarrow\underline\Z_{U_r}
\longrightarrow0
\]

is a projective resolution. Replacing each free integral sheaf by
its \(R\)-module counterpart gives a projective resolution
\(P_\bullet^R\to\underline R_{U_r}\).\leanmark{resolution}
\end{proposition}

\begin{proof}
All terms are projective by \cref{lem:cells,lem:projective-opens}.
For exactness, take a stalk at \(x\in U_r\). The stalk of \(P(W)\)
is zero when \(x\notin W\); if \(x\in V\subseteq W\), the map
\(P(V)_x\to P(W)_x\) is an isomorphism. These assertions follow
from the representing property, or from the stalk--skyscraper
adjunction.

Let \(a\) be the least index with \(x\in W_a\). The surviving
faces form the simplex of charts containing \(x\), and its
augmented chain complex contracts to the vertex \(a\). Explicitly,
prepend \(a\) to an increasing face whose first vertex is larger
than \(a\), using the inverse of the stalk inclusion; on a face
starting at \(a\), set the contraction to zero. Faces starting below
\(a\) have zero stalk. In the alternating boundary formula, deleting
the inserted vertex gives the original face and all other terms
cancel in pairs. Thus \(dh+hd=1\), with the augmentation section
included in degree zero. The augmented stalk complex is exact.
Exactness of sheaves is detected on stalks. The same contraction works
for the module sheaves.
\end{proof}

Applying \(\Hom(-,\mathcal M)\) now computes ordinary sheaf
cohomology, since

\[
H^q(U_r;\mathcal M)
=\Ext^q_{\Sh(U_r;\Ab)}(\underline\Z,\mathcal M).
\]

This also identifies the cochain complex explicitly: its degree-\(k\)
term is the product of the section groups on intersections of \(k+1\)
charts, and its differential is the alternating sum of restrictions.
In particular, with constant \(\F\) coefficients, these sections are
locally constant \(\F\)-valued functions. Because there are no faces
of degree greater than \(r\), we have proved

\[
H^q(U_r;\mathcal M)=0\qquad(q>r)
\]

for every abelian coefficient sheaf.\leanmark{cohomology}

The remaining question is whether a particular cochain in the last
possible degree is a boundary. We first identify that cochain as a
power of a degree-one class.
